\documentclass[10pt,a4paper]{amsart}
\usepackage[margin=19mm]{geometry}
\usepackage{amsmath,amssymb,mathtools}
\usepackage[scr=rsfs,cal=euler]{mathalfa}
\usepackage{xfrac}
\usepackage[unicode,hidelinks,bookmarks=false]{hyperref}
\newcommand{\ie}{i.e.\ }
\newcommand{\cf}{cf.\ }
\newcommand{\resp}{respectively\ }
\newcommand{\Subsection}{\subsection}
\providecommand{\nobreakdash}{\nobreak\hbox{-}}
\setlength{\emergencystretch}{2em}
\multlinegap0pt
\usepackage{bm}
\usepackage{bbm}
\usepackage{dsfont}
\usepackage{xspace}
\usepackage{cancel}
\usepackage{mathtools}
\usepackage{amsthm}
\makeatletter
\@ifundefined{lemma}{\newtheorem{lemma}{Lemma}}{}
\makeatother
\newcommand{\re}{\varepsilon}
\renewcommand{\O}{\mathcal{O}}
\renewcommand{\vec}{\boldsymbol}
\title[Asymptotic models for time-domain scattering by small partic]{Asymptotic models for time-domain scattering by small particles of arbitrary shapes}
\author{Maryna Kachanovska, Adrian Savchuk}
\date{Source: arXiv:2511.11103v1 (CC BY 4.0)}
\begin{document}
\maketitle

\noindent\emph{Source and licence.} Asymptotic models for time-domain scattering by small particles of arbitrary shapes. Version arXiv:2511.11103v1 (CC BY 4.0), distributed under the Creative Commons Attribution 4.0 International licence, as recorded in the arXiv licence metadata for this exact version and shown on the abstract page https://arxiv.org/abs/2511.11103v1. Full rights notice: licenses/arxiv-2511.11103-CC-BY-4.0.txt.\par\smallskip

\noindent\textit{Editorial scope.} Counted unit: the Lemma whose source label is \texttt{lifting\_many\_particles\_proof} in the article (source0.tex lines 3004--3018, source label \texttt{lifting\_many\_particles\_proof}), delivered together with the article's own complete original proof of it (source0.tex lines 3019--3081, one \texttt{proof} environment of 63 lines). No mathematics is written, repaired or re-derived: statement and proof are the article's own text line for line. Required context retained uncounted: lemma:lifting (lines 2942--2956). Every definition, display and label of those lines is retained unchanged. Omitted from this package and not redistributed: 1--2941, 2957--3003, 3082--3312. Nothing inside a retained excerpt is omitted or altered: every retained body is delivered exactly as the article's lines are written, so no omission receipt of any kind accompanies this package. Citations: the retained excerpts carry no citation command at all, so no bibliography entry of the article is transcribed and no reference is redistributed. Reference adaptation: none was needed and none was made; every reference inside the delivered unit resolves inside it, so no unresolved reference or citation is produced. Printed slips kept literal: no printed word, symbol, hypothesis or number of the counted statement or of its proof was altered. Layout and class: the article's own class is \texttt{article} with packages including amsmath, amsthm, amssymb, soul, xcolor, geometry, cite, dsfont, appendix, graphicx, esvect, booktabs, hyperref, bbm; the wrapper is the lane's pinned \texttt{amsart} editorial wrapper at 10pt on A4, which restates the theorem-like environments the retained text needs and adds the article's own macro definitions that the retained text needs (\texttt{\textbackslash O}, \texttt{\textbackslash re}, \texttt{\textbackslash vec}), with extra wrapper packages (bm, bbm, dsfont, xspace, cancel, mathtools). The delivered numbering is the unit's own. Assets: no retained span contains a figure environment, a graphics-inclusion command, a PDF-page inclusion, a table or a TikZ picture; no image file is copied into this package, the package declares no asset dependency and no image file is redistributed with it. Licence: version arXiv:2511.11103v1 of this article is distributed under the Creative Commons Attribution 4.0 International licence, as recorded in the arXiv licence metadata for this exact version and shown on the abstract page https://arxiv.org/abs/2511.11103v1. A full rights notice is delivered at licenses/arxiv-2511.11103-CC-BY-4.0.txt. Characters: every retained excerpt is pure ASCII, so the delivered engine, XeLaTeX with its default Latin Modern text font, reports no missing character.
		\begin{lemma}[Lifting lemma for the single-particle case]
			\label{lemma:lifting}
			Let $a > 0 $, $g\in H^{1/2}(\partial\O^{\re})$. Consider, for $0<\re<1$, the family of functions  $(V^{\re})_{\re>0}\subset H^1(\mathbb{R}^3)$ which satisfy the interior boundary-value problem 
			\begin{align*}%
				- \Delta V^{\re} + a^2 V^{\re} = 0~\text{in}~\mathcal{O}^{\re}, \quad \quad \gamma^{-}_0 V^{\re} = g,
			\end{align*}
			and the exterior boundary-value problem 
			\begin{align*}%
				- \Delta V^{\re} + a^2 V^{\re} = 0~\text{in}~\mathcal{O}^{\re,c}, \quad \quad \gamma^{+}_0 V^{\re} = g.
			\end{align*}
			Then there exists $C>0$, independent of $g, \, a, \re$, but depending on $\O$, s.t. for all $a>0$,  $0<\re<1$, the function $V^{\re}$ satisfies the following bound:
			\begin{align*}%
				\| V^{\re} \|^2_{a,\O^{\re}\cup\O^{\re,c}} \leq C \max(1,\re a) \left(\re^{-1} \| \mathbb{P}_0g\|^2_{L^2(\partial \mathcal{O}^{\re})}+ \| \mathbb{P}_{\perp}g\|^2_{H^{1/2}(\partial \mathcal{O}^{\re})}\right).
			\end{align*}
		\end{lemma}

		\begin{lemma}[Lifting lemma for many particles]
			\label{lifting_many_particles_proof}
			Let $a > 0$, $\vec{\lambda}^{\re} \in H^{1/2}(\Gamma^{\re})$ and $\Lambda^{\re} \in H^1(\mathbb{R}^3)$ satisfy the following interior and exterior boundary problems
			\begin{align*}%
				&-\Delta \Lambda^{\re} + a^2 \Lambda^{\re} = 0 ~in ~ \Omega^{\re}, \quad \gamma^-_0 \Lambda^{\re} = \vec{\lambda}^{\re}, \\ 
				&-\Delta \Lambda^{\re} + a^2 \Lambda^{\re} = 0 ~in ~ \Omega^{\re,c}, \quad \gamma^+_0 \Lambda^{\re} = \vec{\lambda}^{\re}.
			\end{align*} 
			Then, 
			
			\begin{align*}%
				\| \Lambda^{\re} \|^2_{a,\mathbb{R}^3} \leq c (\underline{d}^{\re}_*)^{-2} (1+a) \max(1, a^{-2}) \left(\re^{-1}  \| \mathbb{P}_0\vec{\lambda}^{\re} \|^2_{L^2(\Gamma^{\re})}+ \| \mathbb{P}_{\perp}\vec{\lambda}^{\re} \|^2_{H^{1/2}(\Gamma^{\re})}\right),
			\end{align*}
			
			and $c>0$ is a constant independent of $N, \re, d^{\re}_{ij}, a$. 
		\end{lemma}

		\begin{proof}
			The function $\Lambda^{\re}$ defined as above minimizes the squared energy norm $v\mapsto \|v\|^2_{a,\mathbb{R}^3\setminus\Gamma^{\re}}$ for functions $v\in H^1(\mathbb{R}^3\setminus\Gamma^{\re})$, s.t. $\gamma_0 v=\vec{\lambda}^{\re}$. In other words, 
			\[\Lambda^{\re} = \arg\min_{\substack{v \in H^1(\mathbb{R}^3): \\ \gamma_0 v = \boldsymbol{\Lambda}^{\re}}}\| v \|_{a, \mathbb{R}^3}.\]
			If there exists $v^{\re}\in H^1(\mathbb{R}^3)$ s.t. $\gamma_0 v^{\re}=\vec{\lambda}^{\re}$ and  $\|v^{\re}\|_{a,\mathbb{R}^3}\leq C_{\re}\|\vec{\lambda}^{\re}\|_{H^{1/2}(\Gamma^{\re})}$, then, in virtue of the above observation, 
			\begin{equation}
				\label{eq:lambadre}
				\| \Lambda^{\re} \|_{a, \mathbb{R}^3} \leq \| v^{\re} \|_{a,\mathbb{R}^3} \leq C_{\re} \| \vec{\lambda}^{\re} \|_{H^{1/2}(\Gamma^{\re})}.
			\end{equation}
			
			Our goal is now to construct such a function $v$, for which $C_{\re}$ will be easy to compute. The key idea is to define $v$ as a solution of a certain BVP in the vicinity of $\Omega_k^{\re}$, and $0$ otherwise. This relies on an introduction of certain cutoff functions. 
			
			We will need cutoff functions parametrized by two parameters $r, \, d$. First of all, let  
			\begin{align*}%
				\chi_{d}\in C^{\infty}(\mathbb{R}; [0,1]), \quad	\chi_{d}(\rho):=\left\{
				\begin{array}{ll}
					1, & \rho\leq d/8\\
					0, & \rho\geq d/4,\\
					\in [0,1], &\text{ otherwise },
				\end{array}	
				\right.
			\end{align*}
			Evidently, $\chi_d(\rho)=\chi_1(\rho/d)$. Such a function $\chi_1$ can be constructed e.g. as an appropriate mollification of $\mathbbm{1}_{\rho<1/8}$. 
			Next, we define $\chi_{r,d}(\rho):=\chi_d(\rho-r)=\chi_1(\frac{\rho-r}{d})$, $\rho\geq 0$. 
			
			With the above definition, let us introduce a truncation function associated to $\Omega^{\re}_k$:
			\begin{align*}%
				\chi_k^{\re}(\boldsymbol{x}) := \chi_{r^{\re}_k, \underline{d}^{\re}_*} (|\boldsymbol{x} - \boldsymbol{c}_k|), \quad \boldsymbol{x} \in \mathbb{R}^3.
			\end{align*}	 	
			Note that, by construction 
			\begin{align}
				\label{eq:chione}
				\text{supp} \,\chi_k^{\re} &\subseteq B(\boldsymbol{c}_k, r^{\re}_k+\underline{d}^{\re}_*/4),\quad \text{ therefore, }\quad\overline{\Omega^{\re}_k}\subset \operatorname{supp}\chi_k^{\re},
				\quad 
				\left.\chi_k^{\re}\right|_{\overline{\Omega^{\re}_k}}=1,\\
				\label{eq:supchione}
				\text{supp}\,\chi_k^{\re}&  \cap \text{supp}\,\chi_{\ell}^{\re} = \emptyset, \quad \text{if}~k \neq \ell.
			\end{align} 
			Moreover, because $\chi_{r,d}(\rho)=\chi_1\left(\frac{\rho-r}{d}\right)$, with some $C>0$, it holds that  
			\begin{align}
				\label{eq:bounds_on_cut_off}
				\| \chi_k^{\re} \|_{L^{\infty}(\mathbb{R}^3)} \leq 1, \quad \| \nabla \chi_k^{\re} \|_{L^{\infty}(\mathbb{R}^3)}\equiv \|\partial_{\rho}\chi_{r_k^{\re}, \underline{d}^{\re}_*}\|_{L^{\infty}(\mathbb{R}^3)} \leq C (\underline{d}^{\re}_*)^{-1}.	
			\end{align}
			Using the above definition of cut-off functions, let us introduce a lifting $\Lambda^{\re}_{\chi}$ of $\vec{\lambda}^{\re}$ (this will play a role of the function $v^{\re}$ in the beginning of the proof) as follows:
			\begin{align*}%
				v^{\re}:=\Lambda^{\re}_{\chi} := \sum_{k=1}^N \chi_k^{\re} \Lambda^{\re}_k,
			\end{align*}
			where $\Lambda^{\re}_k \in H^{1}(\mathbb{R}^3)$ satisfies the following Dirichlet boundary problem:
			\begin{align*}%
				&- \Delta \Lambda^{\re}_k + a^2  \Lambda^{\re}_k = 0 \text{ in } \Omega^{\re}_k, \quad  \gamma^{-}_0 \Lambda^{\re}_k = \lambda^{\re}_k, \\
				&- \Delta \Lambda^{\re}_k + a^2  \Lambda^{\re}_k = 0 \text{ in } \Omega^{\re, c}_k, \quad  \gamma^{+}_0 \Lambda^{\re}_k = \lambda^{\re}_k.
			\end{align*}
			The properties \eqref{eq:chione} and \eqref{eq:supchione} ensure that  $\gamma_0\Lambda_{\chi}^{\re}=\vec{\lambda}^{\re}$. Let us now estimate the norm of $\Lambda^{\re}_{\chi}$ by definition. In virtue of \eqref{eq:supchione}, $\chi_k\Lambda^{\re}_k$ are mutually orthogonal in $H^1(\mathbb{R}^3)$, and thus 
			\begin{align*}%
				\| \Lambda^{\re}_{\chi} \|^2_{a, \mathbb{R}^3} &=  \|  \sum_{k=1}^N \chi_k^{\re} \Lambda^{\re}_k \|^2_{a,\mathbb{R}^3} = \sum_{k=1}^N \|  \chi_k^{\re} \Lambda^{\re}_k \|^2_{a,\mathbb{R}^3}\\ 
				&=\sum_{k=1}^N \left( a^2 \|  \chi_k \Lambda^{\re}_k \|^2_{L^2(\mathbb{R}^3)}+\|  \nabla\chi_k^{\re} \Lambda^{\re}_k+\chi_k^{\re}\nabla \Lambda^{\re}_k \|^2_{L^2(\mathbb{R}^3)}\right)\\
				&\leq 2\sum_{k=1}^{N}  \left[a^2 \| \chi_k \|^2_{L^{\infty}(\mathbb{R}^3)} \| \Lambda^{\re}_k \|^2_{L^2(\mathbb{R}^3)}  + \| \nabla \chi_k \|^2_{L^{\infty}(\mathbb{R}^3)} \| \Lambda^{\re}_k \|^2_{L^2(\mathbb{R}^3)} + \| \chi_k \|^2_{L^{\infty}(\mathbb{R}^3)} \|\nabla \Lambda^{\re}_k \|^2_{L^2(\mathbb{R}^3)} \right]\\
				&\overset{	\eqref{eq:bounds_on_cut_off}}{\leq }2C\sum_{k=1}^{N} \left[\| \Lambda^{\re}_k \|^2_{a, \mathbb{R}^3} + (\underline{d}^{\re}_*)^{-2} \| \Lambda^{\re}_k \|^2_{L^2(\mathbb{R}^3)}\right] \\
				&\leq 2C \sum_{k=1}^{N} \left[\| \Lambda^{\re}_k \|^2_{a, \mathbb{R}^3} + ( \underline{d}^{\re}_*)^{-2} a^{-2} \| \Lambda^{\re}_k \|^2_{a, \mathbb{R}^3}\right] \\
				&\leq 2C(1+ a^{-2}(\underline{d}^{\re}_*)^{-2}) \sum_{k=1}^{N} \| \Lambda^{\re}_k \|^2_{a, \mathbb{R}^3}\leq 2C(1+ a^{-2})\max(1, (\underline{d}^{\re}_*)^{-2})\sum_{k=1}^{N} \| \Lambda^{\re}_k \|^2_{a, \mathbb{R}^3}\\
				&= 4C\max(1,a^{-2}) (\underline{d}^{\re}_*)^{-2}\sum_{k=1}^{N} \| \Lambda^{\re}_k \|^2_{a, \mathbb{R}^3}.
			\end{align*}
			Next, we apply Lemma \ref{lemma:lifting} to bound each $\Lambda_k^{\re}$ in terms of $\lambda^{\re}_k$ with the bound $\max(1, \re a) \leq (1+a)$, and the desired result follows at once by \eqref{eq:lambadre} with $v^{\re}=\Lambda^{\re}_{\chi}$.
		\end{proof}

\end{document}
